% New English translation from checked cote 25, site 75 (PDF page 76).
% AI-translated by GPT-6 Astra (OpenAI, ChatGPT Pro mode).
\scanpage{25/75}
\mspage{70}

20.13. \uline{Elementary morphisms and the theorem of} M. \uline{Artin}.

\uline{Definition} 13.1. A morphism $f:X\hand{\to} Y$ of preschemes is called an “\var{\xout{elem}}{}{typewriter strike-out}{ega25_25_075-v0001}
elementary morphism” if X is \var{\xout{a}}{}{typewriter strike-out}{ega25_25_075-v0002} Y-isomorphic to a prescheme of the form
$X' - Z$, where $X'\hand{\to} X$ is a smooth projective morphism with connected geometric
fibers\var{\xout{\unsure{,}}}{}{typewriter strike-out}{ega25_25_075-v0003} of dimension 1, and where Z is a closed subprescheme
of X', \var{\xout{étale over}}{}{typewriter strike-out}{ega25_25_075-v0004} such that the morphism $Z\hand{\to} Y$ is étale, surjective, and
of constant degree. A morphism is called \var{\xout{\unsure{multip-elementary}}}{}{typewriter strike-out}{ega25_25_075-v0005} poly-elementary
if it is a composite of elementary morphisms. A prescheme X
over a field k is called poly-elementary (over k) if the structural
morphism $X\hand{\to} \mathrm{Spec}(k)$ is poly-elementary.

\uline{Theorem} 13.2. \var{\inserted[typed]{(M.Artin)}}{(M.Artin)}{insertion}{ega25_25_075-v0006} Let X be a \var{\inserted[typed]{geom. irred.}}{geom. irred.}{insertion}{ega25_25_075-v0007} prescheme over a \var{\inserted[typed]{perfect infinite}}{perfect infinite}{insertion}{ega25_25_075-v0008} field k, x a smooth point
of X; then x has a fundamental system of poly-elementary open
neighborhoods.

Replacing X by a given neighborhood of x, \var{\deleted{one \var{\xout{\unsure{can suppo}}}{}{typewriter strike-out}{ega25_25_075-v0010} es}}{}{deletion}{ega25_25_075-v0009} \var{\inserted{it is}}{it is}{insertion}{ega25_25_075-v0011}
\var{\deleted{reduced to}}{}{deletion}{ega25_25_075-v0012} \var{\inserted{equivalent to}}{equivalent to}{insertion}{ega25_25_075-v0013} showing that there exists an elementary open neighborhood of x in X.
\var{\deleted{We may clearly assume X smooth, and}}{}{deletion}{ega25_25_075-v0014} \var{\inserted[typed]{\var{\deleted{and affine}}{}{deletion}{ega25_25_075-v0016}}}{\var{\deleted{and affine}}{}{deletion}{ega25_25_075-v0016}}{insertion}{ega25_25_075-v0015} \var{\over{r}{R}}{R}{writing over}{ega25_25_075-v0017}easoning by induction on
the dimension n of X, we \var{\xout{may assume}}{}{typewriter strike-out}{ega25_25_075-v0018} are reduced to showing \var{\inserted[typed]{that if $n\hand{>}0$,}}{that if $n\hand{>}0$,}{insertion}{ega25_25_075-v0019} \var{\deleted{that}}{}{deletion}{ega25_25_075-v0020}there exists
an open neighborhood U of x and an elementary morphism $f: U\hand{\to} V$,
V being a smooth scheme over k (necessarily geom. irred \var{\inserted[typed]{and}}{and}{insertion}{ega25_25_075-v0021} of dim n-1).
\var{\xout{When n=1}}{}{typewriter strike-out}{ega25_25_075-v0022} (The case n=0 is clearly trivial, since then
X is k-isomorphic to Spec(k), which is poly-elementary over Spec(k), it being
understood in 13.1 that we do not exc\hand{l}ude the composite of the empty family of
morphisms --- we should have mentioned this in one form or another in
13.1.) The need to assume k perfect beforehand already appears
in the case n=1: we then take for X' the canonical normal projective model
of the function field K of X (cf Chap II) par.7), the fact that
\tnote{25-075-1}{The defining arrow from X-prime to X is retained exactly as in the checked source. The canceled es remains an incomplete fragment. The French correction of que to qu' before alors changes no English word and has no English variant. The handwritten l in exclut is represented by the corresponding letter in exclude.}

\mspage{75}
